\documentclass[a4paper,11pt]{book}
%\documentclass[a4paper,twoside,11pt,titlepage]{book}
\usepackage{listings}
\usepackage[utf8]{inputenc}
\usepackage[spanish]{babel}

% \usepackage[style=list, number=none]{glossary} %
%\usepackage{titlesec}
%\usepackage{pailatino}

\decimalpoint
\usepackage{dcolumn}
\newcolumntype{.}{D{.}{\esperiod}{-1}}
\makeatletter
\addto\shorthandsspanish{\let\esperiod\es@period@code}
\makeatother


%\usepackage[chapter]{algorithm}
\RequirePackage{verbatim}
%\RequirePackage[Glenn]{fncychap}
\usepackage{fancyhdr}
\usepackage{graphicx}
\usepackage{afterpage}

\usepackage{longtable}

\usepackage[pdfborder={000}]{hyperref} %referencia

% ********************************************************************
% Re-usable information
% ********************************************************************
\newcommand{\myTitle}{Título del proyecto\xspace}
\newcommand{\myDegree}{Grado en ...\xspace}
\newcommand{\myName}{Nombre Apllido1 Apellido2 (alumno)\xspace}
\newcommand{\myProf}{Nombre Apllido1 Apellido2 (tutor1)\xspace}
\newcommand{\myOtherProf}{Nombre Apllido1 Apellido2 (tutor2)\xspace}
%\newcommand{\mySupervisor}{Put name here\xspace}
\newcommand{\myFaculty}{Escuela Técnica Superior de Ingenierías Informática y de
Telecomunicación\xspace}
\newcommand{\myFacultyShort}{E.T.S. de Ingenierías Informática y de
Telecomunicación\xspace}
\newcommand{\myDepartment}{Departamento de ...\xspace}
\newcommand{\myUni}{\protect{Universidad de Granada}\xspace}
\newcommand{\myLocation}{Granada\xspace}
\newcommand{\myTime}{\today\xspace}
\newcommand{\myVersion}{Version 0.1\xspace}


\hypersetup{
pdfauthor = {\myName (enrihuescar@correo.ugr.es)},
pdftitle = {\myTitle},
pdfsubject = {},
pdfkeywords = {palabra_clave1, palabra_clave2, palabra_clave3, ...},
pdfcreator = {LaTeX con el paquete ....},
pdfproducer = {pdflatex}
}

%\hyphenation{}


%\usepackage{doxygen/doxygen}
%\usepackage{pdfpages}
\usepackage{url}
\usepackage{colortbl,longtable}
\usepackage[stable]{footmisc}
%\usepackage{index}

%\makeindex
%\usepackage[style=long, cols=2,border=plain,toc=true,number=none]{glossary}
% \makeglossary

% Definición de comandos que me son tiles:
%\renewcommand{\indexname}{Índice alfabético}
%\renewcommand{\glossaryname}{Glosario}

\pagestyle{fancy}
\fancyhf{}
\fancyhead[LO]{\leftmark}
\fancyhead[RE]{\rightmark}
\fancyhead[RO,LE]{\textbf{\thepage}}
\renewcommand{\chaptermark}[1]{\markboth{\textbf{#1}}{}}
\renewcommand{\sectionmark}[1]{\markright{\textbf{\thesection. #1}}}

\setlength{\headheight}{1.5\headheight}

\newcommand{\HRule}{\rule{\linewidth}{0.5mm}}
%Definimos los tipos teorema, ejemplo y definición podremos usar estos tipos
%simplemente poniendo \begin{teorema} \end{teorema} ...
\newtheorem{teorema}{Teorema}[chapter]
\newtheorem{ejemplo}{Ejemplo}[chapter]
\newtheorem{definicion}{Definición}[chapter]

\definecolor{gray97}{gray}{.97}
\definecolor{gray75}{gray}{.75}
\definecolor{gray45}{gray}{.45}
\definecolor{gray30}{gray}{.94}

\lstset{ frame=Ltb,
     framerule=0.5pt,
     aboveskip=0.5cm,
     framextopmargin=3pt,
     framexbottommargin=3pt,
     framexleftmargin=0.1cm,
     framesep=0pt,
     rulesep=.4pt,
     backgroundcolor=\color{gray97},
     rulesepcolor=\color{black},
     %
     stringstyle=\ttfamily,
     showstringspaces = false,
     basicstyle=\scriptsize\ttfamily,
     commentstyle=\color{gray45},
     keywordstyle=\bfseries,
     %
     numbers=left,
     numbersep=6pt,
     numberstyle=\tiny,
     numberfirstline = false,
     breaklines=true,
   }
 
% minimizar fragmentado de listados
\lstnewenvironment{listing}[1][]
   {\lstset{#1}\pagebreak[0]}{\pagebreak[0]}

\lstdefinestyle{CodigoC}
   {
	basicstyle=\scriptsize,
	frame=single,
	language=C,
	numbers=left
   }
\lstdefinestyle{CodigoC++}
   {
	basicstyle=\small,
	frame=single,
	backgroundcolor=\color{gray30},
	language=C++,
	numbers=left
   }

 
\lstdefinestyle{Consola}
   {basicstyle=\scriptsize\bf\ttfamily,
    backgroundcolor=\color{gray30},
    frame=single,
    numbers=none
   }


\newcommand{\bigrule}{\titlerule[0.5mm]}


%Para conseguir que en las páginas en blanco no ponga cabecerass
\makeatletter
\def\clearpage{%
  \ifvmode
    \ifnum \@dbltopnum =\m@ne
      \ifdim \pagetotal <\topskip
        \hbox{}
      \fi
    \fi
  \fi
  \newpage
  \thispagestyle{empty}
  \write\m@ne{}
  \vbox{}
  \penalty -\@Mi
}
\makeatother

\usepackage{pdfpages}
\begin{document}
\begin{titlepage}
 
 
\newlength{\centeroffset}
\setlength{\centeroffset}{-0.5\oddsidemargin}
\addtolength{\centeroffset}{0.5\evensidemargin}
\thispagestyle{empty}

\noindent\hspace*{\centeroffset}\begin{minipage}{\textwidth}

\centering
\includegraphics[width=0.9\textwidth]{imagenes/logo_ugr.jpg}\\[1.4cm]

\textsc{ \Large TRABAJO FIN DE GRADO\\[0.2cm]}
\textsc{ GRADO EN INGENIERÍA INFORMÁTICA}\\[1cm]
% Upper part of the page
% 
% Title
{\Huge\bfseries Desarrollo de Portal Web de Limpieza de Metadatos\\
}
\noindent\rule[-1ex]{\textwidth}{3pt}\\[3.5ex]

\end{minipage}

\vspace{2.5cm}
\noindent\hspace*{\centeroffset}\begin{minipage}{\textwidth}
\centering

\textbf{Autor}\\ {Enrique García González}\\[2.5ex]
\textbf{Directores}\\
{Nombre Apellido1 Apellido2 (tutor1)\\
Nombre Apellido1 Apellido2 (tutor2)}\\[2cm]
\includegraphics[width=0.3\textwidth]{imagenes/etsiit_logo.png}\\[0.1cm]
\textsc{Escuela Técnica Superior de Ingenierías Informática y de Telecomunicación}\\
\textsc{---}\\
Granada, mes de 201
\end{minipage}
%\addtolength{\textwidth}{\centeroffset}
%\vspace{\stretch{2}}
\end{titlepage}



\chapter*{}
%\thispagestyle{empty}
%\cleardoublepage

%\thispagestyle{empty}

%\input{portada/portada_2}



\cleardoublepage
\thispagestyle{empty}

\begin{center}
{\large\bfseries Título del Proyecto: Subtítulo del proyecto}\\
\end{center}
\begin{center}
Enrique García González\\
\end{center}

%\vspace{0.7cm}
\noindent{\textbf{Palabras clave}: Metadatos, Limpieza de DOcumentos, Portal Web, ......}\\

\vspace{0.7cm}
\noindent{\textbf{Resumen}}\\

Poner aquí el resumen.

Debido al auge de internet, cada día son mar archivos los que compartimos en la web, como por
ejemplo fotografías o domumentos PDF. 
\cleardoublepage


\thispagestyle{empty}


\begin{center}
{\large\bfseries Project Title: Project Subtitle}\\
\end{center}
\begin{center}
First name, Family name (student)\\
\end{center}

%\vspace{0.7cm}
\noindent{\textbf{Keywords}: Keyword1, Keyword2, Keyword3, ....}\\

\vspace{0.7cm}
\noindent{\textbf{Abstract}}\\

Write here the abstract in English.

\chapter*{}
\thispagestyle{empty}

\noindent\rule[-1ex]{\textwidth}{2pt}\\[4.5ex]

Yo, \textbf{Enrique García González}, alumno de la titulación TITULACIÓN de la \textbf{Escuela Técnica Superior
de Ingenierías Informática y de Telecomunicación de la Universidad de Granada}, con DNI 14631973V, autorizo la
ubicación de la siguiente copia de mi Trabajo Fin de Grado en la biblioteca del centro para que pueda ser
consultada por las personas que lo deseen.

\vspace{6cm}

\noindent Fdo: Enrique García González

\vspace{2cm}

\begin{flushright}
Granada a X de mes de 201 .
\end{flushright}


\chapter*{}
\thispagestyle{empty}

\noindent\rule[-1ex]{\textwidth}{2pt}\\[4.5ex]

D. \textbf{Nombre Apellido1 Apellido2 (tutor1)}, Profesor del Área de XXXX del Departamento YYYY de la Universidad de Granada.

\vspace{0.5cm}

D. \textbf{Nombre Apellido1 Apellido2 (tutor2)}, Profesor del Área de XXXX del Departamento YYYY de la Universidad de Granada.


\vspace{0.5cm}

\textbf{Informan:}

\vspace{0.5cm}

Que el presente trabajo, titulado \textit{\textbf{Desarrollo de Portal Web de Limpieza de Metadatos}},
ha sido realizado bajo su supervisión por \textbf{Enrique García González}, y autorizamos la defensa de dicho trabajo ante el tribunal
que corresponda.

\vspace{0.5cm}

Y para que conste, expiden y firman el presente informe en Granada a X de mes de 201 .

\vspace{1cm}

\textbf{Los directores:}

\vspace{5cm}

\noindent \textbf{Nombre Apellido1 Apellido2 (tutor1) \ \ \ \ \ Nombre Apellido1 Apellido2 (tutor2)}

\chapter*{Agradecimientos}
\thispagestyle{empty}

       \vspace{1cm}


A mi familia y amigos, por apoyarme en mis decisiones y ayudarme en mis desvarios.



\newpage
%\frontmatter
\tableofcontents
\listoffigures
%S\listoftables
%
%\mainmatter
\setlength{\parskip}{5pt}
\newpage
\chapter{Introducción}


Tras la publicación del Real Decreto 3/2010, del 8 de enero, se regulo el Esquema Nacional de Seguridad en el ámbito 
de la Administración Electrónica. Este documento se creo con la finalidad de crear las condiciones necesarias en el
uso de los medios electrónicos a traves de distintas medidas de seguridad, permitiendo así el correcto funcionamiento de
las Administraciones públicas.


Dentro del esquema, en el Anexo 11: Medidas de seguridad, encontramos un apartado dedicado a la limpieza de documentos
[5.7.6], en el cual se índica que se debe eliminar cualquier información adicional contenida en campos ocultos, 
meta-datos y comentarios o revisiones anteiores, salvo cuando dicha información sea pertinente para el receptor del 
ducumento.

Como se nos indica tambíen en el Esquema, el incumplimiento de esta medida puede perjudicar:

\begin{enumerate}
	\item Al mantenimiento de la confidencialidad de información que no debería haberse revelado al receptor del
	documento.
	\item Al mantenimiento de la confidencialidad de las fuentes u origenes de la información, que no debe 
	conocer el receptor del documento.
	\item A la buena imagen de la organización que difunde el documento por cuanto demuestra un descuido en su buen
	hacer.
\end{enumerate}



\section{Metadatos: Lo que ocultan los documentos}

Cuando hablamos sobre metadatos, nos referimos a aquellos datos de un documento que se refieren a las propiedades del 
mismo. Cada tipo de documento suele contener distintos tipos de metadatos. Un documento ofimático por ejemplo puede 
contener información sobre la fecha de creación y de última modificación, autor, historial de versiones, etc, mientras 
que una fotografía puede contener otros datos como son la orientación o la geolocalización.

La divulgación de esta información puede no ser peligrosa para nuestro sistema, aunque esto no quita que se este
ignorando la confidencialidad de cierta información al compartirla adjunta a los metadatos de los documentos. Ademas, 
esta información puede comprometer otros sectores de seguridad, pudiendose revelar vulnerabilidades del sistema. Un 
ejemplo sería revelar un sofware y su versión del mismo, con lo que la busqueda de vulnerabilidades por un atacante 
podría ser menos costosa.


\section{Responsabilidad sobre los metadatos}

Como se indica en el artículo 12 del Esquema Nacional, la seguridad deberá comprometer a todos los miembros de la 
organización. Es por ello que son los autores de los documentos que se van a compartir por la red y los que los 
comparten responsables de cumplir con la limpieza de los documentos. En el articulo 5 también se dice que la seguridad 
se entendera como un proceso integral constituido todos los elementos técnicos, humanos, materiales y organizativos, 
relacionados con el sistema. Con estos dos articulos también podemos entender que la organización en si también es en 
parte responsable de que se cumpla el Esquema y que es favorable que ponga a disposición recursos para facilitar su 
cumplimiento, como puede ser herramientas para la localización y limpieza de
metadatos.

Por último, cabe destacar que aunque parezca que la limpieza de documentos puede parecer una tarea menos importante que 
otras de las que aparecen en el Esquema, tal como se indica en el documento la debilidad de un sistema la determina su 
punto más frágil, por lo que no se debe descuidar ningún aspecto en cuanto a seguridad se refiere.


\section{Objetivos}

El objetivo de este proyecto es proporcionar una plataforma a la universidad de Granada para que su
personal pueda limpiar de meta-datos los documentos que van a compartir o subir a la red, conforme a lo establecido
en el Esquema Nacional de Seguridad. Este proyecto me fué propuesto por Antonio Muñoz Ropa, jefe de servicio del area
de seguridad del CSIR, mientras realizaba unas practicas de empresa en el area.

\subsection{Objetivos principales}



\section{Referencias}

Esquema nacional de seguirdad

1.1:

http://blog.isecauditors.com/2014/09/los-metadatos-definicion-riesgo-eliminacion.html
%
\chapter{Especificación de Requisitos}

Antes de empezar un proyecto, se debe realizar un analisis de los requisitos, los cuales muestran las condiciones
que necesitamos para conseguir nuestro objetivo.

\section{Requisitos Funcionales}

Los requisitos funcionales marcan como debe actuar el sistema ante las interraciones entre el sistema y los usuarios.
Tienen que reflejar que servicios ofrecerá y como responderá ante ciertas acciones.

\begin{itemize}
	\renewcommand{\labelitemi}{$\bullet$}


	\item \textbf{RF-1.} Gestión de archivos:
	\begin{itemize}
		\item \textbf{RF-1.1} Se debe permitir subir cualquier tipo de archivo.
		\item \textbf{RF-1.2} Si no se pueden extraer los datos de ese tipo de fichero se debe informar de ello.
		\item \textbf{RF-1.3} Los archivos que los usuarios descarguen deberán contener solamente los metadatos
		indicados por ellos
	\end{itemize}

	\item \textbf{RF-2.} Gestión de metadatos:
	\begin{itemize}
		\item \textbf{RF-2.1} Se debe permitir visualizar los metadatos del archivo.
		\item \textbf{RF-2.2} Debe ser posible elegir que metadatos se eliminan.
		\item \textbf{RF-2.3} Los usuarios podrán modificar o añadir nuevos metadatos .
	\end{itemize}

	\item \textbf{RF-3.} Gestión de metadatos:
	\begin{itemize}
		\item \textbf{RF-3.1} Se debe permitir visualizar los metadatos del archivo.
		\item \textbf{RF-3.2} Debe ser posible elegir que metadatos se eliminan.
		\item \textbf{RF-3.3} Los usuarios podrán modificar o añadir nuevos metadatos.
	\end{itemize}


\section{Requisitos no funcionales}

Los requisitos no funcionales describen restricciones y caracteristicas que han de tenerse en cuenta en el desarrollo
del proyecto sin afectar directamente al funcionamiento del mismo

\begin{itemize}
	\item \textbf{RNF-1} Debe realizarse de manera que pueda ejecutarse bajo un servidor web debian/ubuntu
	\item \textbf{RNF-2} La implementación se realizara en Python usando Django como Framework.
	\item \textbf{RNF-3} Cada usuario solo podrá trabajar sobre un único archivo al mismo tiempo.
	\item \textbf{RNF-4} Un usuario no podra acceder a los archivos de otro usuario.
	\item \textbf{RNF-5} Debe realizarse utilizando solo herramientas de sofware libre.
	\item \textbf{RNF-6} 

\end{itemize}

\end{itemize}
%
%\input{capitulos/03_Planificacion}
%
%\input{capitulos/04_Analisis}
%
%\input{capitulos/05_Diseno}
%
%\input{capitulos/06_Implementacion}
%
%\input{capitulos/07_Pruebas}
%
%\input{capitulos/08_Conclusiones}
%
%%\chapter{Conclusiones y Trabajos Futuros}
%
%
%%\nocite{*}
%\bibliography{bibliografia/bibliografia}\addcontentsline{toc}{chapter}{Bibliografía}
%\bibliographystyle{miunsrturl}
%
%\appendix
%\input{apendices/manual_usuario/manual_usuario}
%%\input{apendices/paper/paper}
%\input{glosario/entradas_glosario}
% \addcontentsline{toc}{chapter}{Glosario}
% \printglossary
\chapter*{}
\thispagestyle{empty}

\end{document}
